El novembre del 2017, diversos observatoris de tot Europa van enregistrar un augment de les concentracions de l'isòtop ruteni 106 (${}^{106}\mathrm{Ru}$). Es desconeix la naturalesa exacta de l'accident que va provocar aquesta emissió radioactiva. Es calcula que, en el moment de l'emissió, l'activitat de la fuita era de $200\ \mathrm{TBq}$ aproximadament. El període de semidesintegració d'aquest isòtop és de 374~dies i es transforma en rodi 106 (${}^{106}\mathrm{Rh}$).

\begin{apartat}{1}
Escriviu l'equació nuclear de la desintegració del ruteni 106, incloent-hi tots els subíndexs i superíndexs, així com els noms de totes les partícules que intervenen en l'equació. Com s'anomena aquesta desintegració?
\end{apartat}

\begin{apartat}{1}
Calculeu l'activitat del ruteni 106 al cap de set mesos (210 dies) de ser alliberat a causa de l'accident.
\end{apartat}

\dades{%
  Nombre atòmic del ruteni, $Z(\mathrm{Ru}) = 44$.\\
  Nombre atòmic del rodi, $Z(\mathrm{Rh}) = 45$.}
